\chapter{Unsupervised Learning and Recommender Systems for DSS}
\label{ch:unsupervised-learning-and-recommender-systems-for-dss}

\section{Chapter Overview and Learning Objectives}

Many decision-support problems do not begin with labeled outcomes. Instead,
they begin with raw operational, textual, or behavioral data where the system
must first discover structure before it can recommend an action. Unsupervised
learning methods help DSS identify segments, anomalies, latent topics, and
behavioral similarities, while recommender systems translate those patterns
into ranked options for users or managers.

\subsection*{Learning objectives}
\addcontentsline{toc}{subsection}{Learning objectives}
By the end of this chapter, the student should be able to:
\begin{itemize}
  \item distinguish unsupervised learning from supervised learning in DSS,
  \item explain how clustering, dimensionality reduction, and association
  analysis support decision making,
  \item describe the architecture of a recommender system in a DSS context,
  \item compare content-based, collaborative, and hybrid recommendation
  strategies,
  \item identify evaluation criteria for unsupervised and recommender systems,
  \item discuss risks such as popularity bias, cold start, and feedback loops.
\end{itemize}

\section{Unsupervised Learning in DSS}

\subsection{Why unlabeled structure matters}
In practice, organizations often have abundant data but few reliable labels.
Unsupervised learning helps DSS answer questions such as:
\begin{itemize}
  \item Which customers or students form meaningful segments?
  \item Which records look unusual enough to review?
  \item Which variables summarize the main structure of the data?
\end{itemize}

\subsection{Clustering for segmentation}
Clustering groups similar records without predefined classes. In DSS,
clustering is often used for:
\begin{itemize}
  \item customer segmentation for differentiated interventions,
  \item patient grouping for care pathways,
  \item student grouping for learning support,
  \item incident grouping for operational triage.
\end{itemize}
Common methods include $k$-means, hierarchical clustering, and density-based
methods such as DBSCAN.

\subsection{Dimensionality reduction}
High-dimensional data can be difficult to interpret and visualize.
Dimensionality reduction methods such as Principal Component Analysis (PCA) and
manifold-learning techniques help summarize the dominant structure of the data.
In DSS, this can support exploratory analysis, dashboard simplification, and
feature compression before downstream modeling.

\subsection{Association analysis}
Association rules discover co-occurrence patterns such as items frequently
purchased together, services used in sequence, or symptoms commonly observed
together. While association analysis does not by itself prescribe action, it
can drive bundling, cross-selling, fraud review, and knowledge discovery.

\section{From Patterns to Decisions}

\subsection{Anomaly and novelty detection}
Anomaly detection is useful when the objective is to identify cases that merit
review rather than produce a final automated decision. Examples include:
\begin{itemize}
  \item suspicious transactions,
  \item unusual patient deterioration signals,
  \item abnormal network events,
  \item unexpected operational bottlenecks.
\end{itemize}
The DSS should treat anomaly scores as prioritization signals and pair them
with human review workflows.

\subsection{Interpretation challenges}
Unsupervised models can produce patterns that are mathematically coherent but
operationally meaningless. For this reason, domain validation is essential:
\begin{itemize}
  \item inspect cluster profiles,
  \item test stability across time and samples,
  \item confirm that discovered patterns support actual decisions.
\end{itemize}

\section{Recommender Systems}

\subsection{What a recommender system does}
A recommender system ranks items, actions, or interventions for a user, group,
or case. In DSS, recommendation may target:
\begin{itemize}
  \item products or services,
  \item learning resources,
  \item treatment pathways,
  \item support actions for at-risk cases.
\end{itemize}

\subsection{Content-based recommendation}
Content-based recommenders rely on attributes of items and profiles of users or
cases. They are useful when rich metadata is available and when explaining the
recommendation in terms of item attributes is important.

\subsection{Collaborative filtering}
Collaborative filtering uses similarity in historical behavior across users and
items. It can reveal patterns not obvious from metadata alone, but it suffers
from cold-start problems when new users or new items have little interaction
history.

\subsection{Hybrid recommendation}
Many production DSS use hybrid recommenders that combine content features,
interaction history, business rules, and constraints. Hybrid systems are often
more robust and better aligned with policy requirements.

\section{Evaluation and Governance}

\subsection{Evaluating unsupervised methods}
Unsupervised methods are evaluated differently from supervised models. Typical
criteria include:
\begin{itemize}
  \item internal structure metrics such as silhouette score,
  \item stability across runs or time periods,
  \item interpretability and usefulness to decision makers,
  \item downstream impact on decisions and interventions.
\end{itemize}

\subsection{Evaluating recommenders}
Recommendation quality can be assessed using:
\begin{itemize}
  \item ranking metrics such as precision at $k$ and recall at $k$,
  \item online response metrics such as click-through or acceptance rate,
  \item diversity and novelty of recommendations,
  \item fairness and exposure across groups or items.
\end{itemize}

\subsection{Risks and limitations}
Important governance issues include:
\begin{itemize}
  \item \textbf{Cold start}: insufficient history for new users or items.
  \item \textbf{Popularity bias}: over-exposure of already popular items.
  \item \textbf{Feedback loops}: recommendations shape future data.
  \item \textbf{Filter bubbles}: excessive narrowing of user options.
  \item \textbf{Opaque segments}: clusters that are hard to justify in practice.
\end{itemize}

\section{Summary and Concluding Remarks}

This chapter showed how unlabeled data can still support decision making
through clustering, dimensionality reduction, association analysis, anomaly
detection, and recommendation. The main design principle is to connect
discovered structure to concrete actions: segmentation must lead to
interventions, anomaly scores must feed review queues, and recommendations must
respect organizational constraints and fairness requirements.

\section{Key Terms}
\begin{description}[style=nextline]
  \item[Unsupervised learning] Learning patterns from unlabeled data.
  \item[Clustering] Grouping similar observations without predefined labels.
  \item[Dimensionality reduction] Compressing high-dimensional data into a
  smaller representation while preserving important structure.
  \item[Association rule] A pattern describing items or events that frequently
  occur together.
  \item[Anomaly detection] Identifying unusual records or behaviors that may
  require review.
  \item[Recommender system] A system that ranks items or actions for a user or
  case.
  \item[Collaborative filtering] Recommendation using shared behavioral
  patterns across users and items.
  \item[Content-based recommendation] Recommendation using item and user
  attributes.
  \item[Cold start] The difficulty of recommending for new users or items with
  limited history.
  \item[Popularity bias] Over-recommending already popular items.
\end{description}

\section{Multiple-Choice Questions (MCQs)}

\subsection*{Questions}
\addcontentsline{toc}{subsection}{Questions}
\begin{enumerate}[label=\textbf{Q\arabic*.}]
  \item Clustering is most appropriate when:
  \begin{enumerate}[label=\alph*)]
    \item class labels are already known
    \item the DSS needs to discover latent segments in unlabeled data
    \item the objective is to solve a linear program
    \item only deterministic rules are allowed
  \end{enumerate}

  \item A common cold-start problem occurs when:
  \begin{enumerate}[label=\alph*)]
    \item there is too much labeled data
    \item a recommender has little history for new users or items
    \item clusters overlap perfectly
    \item dimensionality is low
  \end{enumerate}

  \item Which issue is most associated with recommender-system feedback loops?
  \begin{enumerate}[label=\alph*)]
    \item recommendations change future behavior and future data
    \item the system cannot rank items
    \item rule-based reasoning becomes impossible
    \item all items receive equal exposure
  \end{enumerate}

  \item PCA is primarily used for:
  \begin{enumerate}[label=\alph*)]
    \item human approval workflows
    \item dimensionality reduction and structure summarization
    \item legal compliance audits
    \item backward chaining
  \end{enumerate}
\end{enumerate}

\subsection*{Answer Key}
\addcontentsline{toc}{subsection}{Answer Key}
\begin{enumerate}[label=\textbf{Q\arabic*.}]
  \item b
  \item b
  \item a
  \item b
\end{enumerate}
